\section{AHSME 1978}
        \begin{enumerate}
        	\item (\textit{AHSME 1978}) If $1-\frac{4}{x}+\frac{4}{x^2}=0$, then $\frac{2}{x}$ equals


 $\textbf{(A) }-1\qquad \textbf{(B) }1\qquad \textbf{(C) }2\qquad \textbf{(D) }-1\text{ or }2\qquad  \textbf{(E) }-1\text{ or }-2$

 \textbf{Answer: }B



\textbf{Solution 1}

By guessing and checking, 2 works. $\frac{2}{x} = \boxed{\textbf{(B)  }1}$~awin

 



\textbf{Solution 2}

Multiplying each side by $x^2$, we get $x^2-4x+4 = 0$. Factoring, we get $(x-2)(x-2) = 0$. Therefore, $x = 2$. $\frac{2}{x} = \boxed{\textbf{(B)  }1}$~awin

 



\textbf{Solution 3}

Directly factoring, we get $(1-\frac{2}{x})^2 = 0$. Thus $\frac{2}{x}$ must equal $\boxed{\textbf{(B)  }1}$



	\item (\textit{AHSME 1978}) If four times the reciprocal of the circumference of a circle equals the diameter of the circle, then the area of the circle is 


 $\textbf{(A) }\frac{1}{\pi^2}\qquad \textbf{(B) }\frac{1}{\pi}\qquad \textbf{(C) }1\qquad \textbf{(D) }\pi\qquad \textbf{(E) }\pi^2$

 \textbf{Answer: }C



\textbf{Solution 1}

Creating equations, we get $4\cdot\frac{1}{2\pi r} = 2r$. Simplifying, we get $\frac{1}{\pi r} = r$. Multiplying each side by $r$, we get $\frac{1}{\pi} = r^2$. Because the formula of the area of a circle is $\pi r^2$, we multiply each side by $\pi$ to get $1 = \pi r^2$.Therefore, our answer is $\boxed{\textbf{(C)  }1}$



	\item (\textit{AHSME 1978}) For all non-zero numbers $x$ and $y$ such that $x = 1/y$, $\left(x-\frac{1}{x}\right)\left(y+\frac{1}{y}\right)$ equals


 $\textbf{(A) }2x^2\qquad \textbf{(B) }2y^2\qquad \textbf{(C) }x^2+y^2\qquad \textbf{(D) }x^2-y^2\qquad  \textbf{(E) }y^2-x^2$

 \textbf{Answer: }D



\textbf{Solution 1}

Using substitution, we can substitute y into the equation in the first parentheses. Therefore, we'll get 
 \[\left(x-\frac{1}{\frac{1}{y}}\right)\left(y+\frac{1}{y}\right)\]

 Because $x = \frac{1}{y}$, we can also see that $y = \frac{1}{x}$. Using substitution again, we can substitute x into the second equation getting
 \[\left(x-\frac{1}{\frac{1}{y}}\right)\left(y+\frac{1}{\frac{1}{x}}\right)\]

 Simplifying, we get $(x-y)(y+x)$. Multiplying, we get $\boxed{\textbf{(D)   }x^2-y^2}$



	\item (\textit{AHSME 1978}) If $a = 1,~ b = 10, ~c = 100$, and $d = 1000$, then $(a+ b+ c-d) + (a + b- c+ d) +(a-b+ c+d)+ (-a+ b+c+d)$ is equal to


 $\textbf{(A) }1111\qquad \textbf{(B) }2222\qquad \textbf{(C) }3333\qquad \textbf{(D) }1212\qquad  \textbf{(E) }4242$

 \textbf{Answer: }B



\textbf{Solution 1}

Adding all four of the terms up, we can see that it equals 
 \[2(a+b+c+d)\]This is equal to $2(1111) = \boxed{\textbf{(B) }2222}$ ~awin

 


	\item (\textit{AHSME 1978}) Four boys bought a boat for $\textdollar 60$. The first boy paid one half of the sum of the amounts paid by the other boys; 
the second boy paid one third of the sum of the amounts paid by the other boys; 
and the third boy paid one fourth of the sum of the amounts paid by the other boys. How much did the fourth boy pay? 


 $\textbf{(A) }\textdollar 10\qquad \textbf{(B) }\textdollar 12\qquad \textbf{(C) }\textdollar 13\qquad \textbf{(D) }\textdollar 14\qquad \textbf{(E) }\textdollar 15$

 \textbf{Answer: }C



\textbf{Solution 1}

If the first boy paid one half of the sum of the amounts paid by the other boys, that means that he paid $\frac{1}{3}$ of the total.If the second boy paid one third of the sum of the amounts paid by the other boys, that means that he paid $\frac{1}{4}$ of the total.If the third boy paid one fourth of the sum of the  amounts paid by the other boys, that means that he paid $\frac{1}{5}$ of the total. 

 Summing it up, we get $\textdollar 20 + \textdollar 15 + \textdollar 12 = \textdollar 47$.Therefore, our answer is $\textdollar 60 - \textdollar 47 = \boxed{\textbf{(C) }\textdollar 13}$~awin

 


	\item (\textit{AHSME 1978}) The number of distinct pairs $(x,y)$ of real numbers satisfying both of the following equations: 


 
 \[x=x^2+y^2 \    \ y=2xy\] 
is


 $\textbf{(A) }0\qquad \textbf{(B) }1\qquad \textbf{(C) }2\qquad \textbf{(D) }3\qquad  \textbf{(E) }4$

 \textbf{Answer: }E



\textbf{Solution}

If $x=x^2+y^2$ and $y=2xy$, then we can break this into two cases.

 Case 1: $y = 0$

 If $y = 0$, then $x = x^2$ and $0 = 0$

 Therefore, $x = 0$ or $x = 1$

 This yields 2 solutions

 Case 2: $x = \frac{1}{2}$ 

 If $x = \frac{1}{2}$, this means that $y = y$, and $\frac{1}{2} = \frac{1}{4} + y^2$. 

 Because y can be negative or positive, this yields $y = \frac{1}{2}$ or $y = -\frac{1}{2}$

 This yields another 2 solutions.

 $2+2 = \boxed{\textbf{(E) 4}}$



	\item (\textit{AHSME 1978}) Opposite sides of a regular hexagon are $12$ inches apart. The length of each side, in inches, is


 $\textbf{(A) }7.5\qquad \textbf{(B) }6\sqrt{2}\qquad \textbf{(C) }5\sqrt{2}\qquad \textbf{(D) }\frac{9}{2}\sqrt{3}\qquad \textbf{(E) }4\sqrt{3}$

 \textbf{Answer: }E



\textbf{Solution}

Draw a perpendicular through the midpoint of the line of length $12$ such that it passes through a vertex. We now have created $2$ $30-60-90$ triangles. Using the ratios, we get that the hypotenuse is $6 \times \frac {2}{\sqrt{3}}$ $= \frac {12}{\sqrt {3}}$ $= 4\sqrt{3} \rightarrow \boxed {\textbf{(E)}}$



	\item (\textit{AHSME 1978}) If $x\neq y$ and the sequences $x,a_1,a_2,y$ and $x,b_1,b_2,b_3,y$ each are in arithmetic progression, then $(a_2-a_1)/(b_2-b_1)$ equals


 $\textbf{(A) }\frac{2}{3}\qquad \textbf{(B) }\frac{3}{4}\qquad \textbf{(C) }1\qquad \textbf{(D) }\frac{4}{3}\qquad \textbf{(E) }\frac{3}{2}$

 \textbf{Answer: }D



\textbf{Solution}

WLOG, let $x =2$ and $y = 5$. From the first sequence, we get
 \[2, 3, 4, 5\] so $a_2 - a_1 = 1$From the second sequence, we get
 \[2, 2+r, 2+2r, 2+3r, 5\] so $2+4r = 5$ and $r = \frac{3}{4}$Thus, we have 
 \[2, 2 \frac{3}{4}, 3 \frac{1}{2} ...\] and $b_2 - b_1 = \frac{3}{4}$So $\frac{1}{\frac{3}{4}} = \frac{4}{3}$
 \[\boxed{D}\]



	\item (\textit{AHSME 1978}) If $x<0$, then $\left|x-\sqrt{(x-1)^2}\right|$ equals


 $\textbf{(A) }1\qquad \textbf{(B) }1-2x\qquad \textbf{(C) }-2x-1\qquad \textbf{(D) }1+2x\qquad  \textbf{(E) }2x-1$

 \textbf{Answer: }B



\textbf{Solution}

We have $\sqrt{x^2} = |x|$, so we rewrite the expression as follows.
 \[|x - \sqrt{(x-1)^2}| = |x - |x-1||\]We know that $x < 0$, so $x-1 < 0$. Thus, we can rewrite $|x-1|$ as $1-x$. So
 \[|x - |x-1|| = |x - (1-x)| = |2x - 1|\].Since $x< 0, 2x-1 < 0$. Thus, we can write this as
 \[|2x - 1| = 1- 2x\]$\boxed{B}$



	\item (\textit{AHSME 1978}) If $\mathit{B}$ is a point on circle $\mathit{C}$ with center $\mathit{P}$, then the set of all points $\mathit{A}$ in the plane of circle 
$\mathit{C}$ such that the distance between $\mathit{A}$ and $\mathit{B}$ is less than or equal to the distance between $\mathit{A}$ 
and any other point on circle $\mathit{C}$ is


 $\textbf{(A) }\text{the line segment from }P \text{ to }B\qquad\\ \textbf{(B) }\text{the ray beginning at }P \text{ and passing through }B\qquad\\ \textbf{(C) }\text{a ray beginning at }B\qquad\\ \textbf{(D) }\text{a circle whose center is }P\qquad\\ \textbf{(E) }\text{a circle whose center is }B$

 \textbf{Answer: }B



\textbf{Solution}

Begin by drawing circle P and point B. To satisfy the conditions of the problem, A needs to be in a position where it is closer to B. This can only happen if A and B are on the same line, so we choose from answer choices A and B. 

 We can pick some arbitrary point A outside circle P that is collinear with B and see that the conditions still hold, so the answer is $\boxed{\textbf{(B)}}.$



	\item (\textit{AHSME 1978}) If $r$ is positive and the line whose equation is $x + y = r$ is tangent to the circle whose equation is $x^2 + y ^2 = r$, then $r$ equals


 $\textbf{(A) }\frac{1}{2}\qquad \textbf{(B) }1\qquad \textbf{(C) }2\qquad \textbf{(D) }\sqrt{2}\qquad  \textbf{(E) }2\sqrt{2}$

 \textbf{Answer: }C



\textbf{Solution}

The circle $x^2 + y^2 = r$ has center $(0,0)$ and radius $\sqrt{r}$. Therefore, if the line $x + y = r$ is tangent to the circle $x^2 + y^2 = r$, then the distance between $(0,0)$ and the line $x + y = r$ is $\sqrt{r}$.

 The distance between $(0,0)$ and the line $x + y = r$ is
 \[\frac{|0 + 0 - r|}{\sqrt{1^2 + 1^2}} = \frac{r}{\sqrt{2}}.\]Hence,
 \[\frac{r}{\sqrt{2}} = \sqrt{r}.\]Then $r = \sqrt{r} \cdot \sqrt{2}$, so $\sqrt{r} = \sqrt{2}$, which means $r = \boxed{2}$ or (C), $2$.

 


	\item (\textit{AHSME 1978}) In $\triangle ADE$, $\measuredangle ADE=140^\circ$, points $B$ and $C$ lie on sides $AD$ and $AE$, 
respectively, and points $A,~B,~C,~D,~E$ are distinct.* If lengths $AB,~BC,~CD$, and $DE$ are all equal, 
then the measure of $\measuredangle EAD$ is


 
            \begin{itemize}
        \item The specification that points $A,B,C,D,E$ be distinct was not included in the original statement of the problem.
            \end{itemize}
            If $B=D$, then $C=E$ and $\measuredangle EAD=20^\circ$.     


 $\textbf{(A) }5^\circ\qquad \textbf{(B) }6^\circ\qquad \textbf{(C) }7.5^\circ\qquad \textbf{(D) }8^\circ\qquad \textbf{(E) }10^\circ$

 \textbf{Answer: }E


	\item (\textit{AHSME 1978}) If $a,b,c$, and $d$ are non-zero numbers such that $c$ and $d$ are the solutions of $x^2+ax+b=0$ and $a$ and $b$ are 
the solutions of $x^2+cx+d=0$, then $a+b+c+d$ equals


 $\textbf{(A) }0\qquad \textbf{(B) }-2\qquad \textbf{(C) }2\qquad \textbf{(D) }4\qquad  \textbf{(E) }(-1+\sqrt{5})/2$

 \textbf{Answer: }B



\textbf{Solution}

By Vieta's formulas, $c + d = -a$, $cd = b$, $a + b = -c$, and $ab = d$. From the equation $c + d = -a$, $d = -a - c$, and from the equation $a + b = -c$, $b = -a - c$, so $b = d$.

 Then from the equation $cd = b$, $cb = b$. Since $b$ is nonzero, we can divide both sides of the equation by $b$ to get $c = 1$. Similarly, from the equation $ab = d$, $ab = b$, so $a = 1$. Then $b = d = -a - c = -2$. Therefore, $a + b + c + d = 1 + (-2) + 1 + (-2) = \boxed{-2}$. The answer is (B).

 


	\item (\textit{AHSME 1978}) If an integer $n > 8$ is a solution of the equation $x^2 - ax+b=0$ and the representation of $a$ in the base-$n$ number system is $18$, 
then the base-n representation of $b$ is


 $\textbf{(A)}\ 18 \qquad \textbf{(B)}\ 20 \qquad \textbf{(C)}\ 80 \qquad \textbf{(D)}\ 81 \qquad \textbf{(E)}\ 280$

 \textbf{Answer: }C



\textbf{Solution}

Assuming the solutions to the equation are n and m, by Vieta's formulas, $n_n + m_n = 18_n$.  

 $n_n = 10_n$, so $10_n + m_n = 18_n$.  

 
 \[m_n = 8_n\].

 Also by Vieta's formulas, $n_n \cdot m_n = b_n$. 
 \[10_n \cdot 8_n = \boxed{80_n}\].

 The answer is (C) $80$



	\item (\textit{AHSME 1978}) If $\sin x+\cos x=1/5$ and $0\le x<\pi$, then $\tan x$ is


 $\textbf{(A) }-\frac{4}{3}\qquad \textbf{(B) }-\frac{3}{4}\qquad \textbf{(C) }\frac{3}{4}\qquad \textbf{(D) }\frac{4}{3}\qquad\\ \textbf{(E) }\text{not completely determined by the given information}$

 \textbf{Answer: }A



\textbf{Solution}

Squaring the equation, we get
 \[\sin ^2 x + \cos ^2 x + 2 \sin x \cos x = \frac{1}{25} =\Rrightarrow 2\sin x \cos x = -\frac{24}{25} \Rrightarrow \sin x \cos x = -\frac{12}{25}\]Recall that $(a-b)^2 = (a+b)^2 - 4ab$, so
 \[(\sin x - \cos x)^2 = (\sin x + \cos x)^2 - 4 \sin x \cos x\]
 \[(\sin x - \cos x)^2 = \frac{1}{25} - 4 (- \frac{12}{25})\]
 \[(\sin x - \cos x) = \frac{7}{5}\]We can now solve for the values of $\sin x$ and $\cos x$
 \[\sin x + \cos x = \frac{1}{5}\]
 \[\sin x - \cos x = \frac{7}{5}\]
 \[2 \sin x = \frac{8}{5} \Rrightarrow \sin x = \frac{4}{5}, \cos x = -\frac{3}{5}\]Since $\tan x = \frac{\sin x}{\cos x}$, we have
 \[\tan x = \frac{\frac{4}{5}}{-\frac{3}{5}} = -\frac{4}{3}\]$\boxed{A}$



	\item (\textit{AHSME 1978}) In a room containing $N$ people, $N > 3$, at least one person has not shaken hands with everyone else in the room. 
What is the maximum number of people in the room that could have shaken hands with everyone else?


 $\textbf{(A) }0\qquad \textbf{(B) }1\qquad \textbf{(C) }N-1\qquad \textbf{(D) }N\qquad  \textbf{(E) }\text{none of these}$

 \textbf{Answer: }E



\textbf{Solution}

We assume that one person hasn't shaken hands with the other N people, meaning that they only had shaken hands with N-1 people. However, this doesn't make sense! There will be 1 person ( from the N-1) people who hasn't shaken hands with the first person, meaning that they also only had shaken hands with N-1 people. Therefore, there is a minimum of 2 people in the room that haven't shaken hands with one person. Therefore, the maximum is N-2 people, so the answer is $E$



	\item (\textit{AHSME 1978}) If $k$ is a positive number and $f$ is a function such that, for every positive number $x$, $\left[f(x^2+1)\right]^{\sqrt{x}}=k$; 
then, for every positive number $y$, $\left[f(\frac{9+y^2}{y^2})\right]^{\sqrt{\frac{12}{y}}}$ is equal to


 $\textbf{(A) }\sqrt{k}\qquad \textbf{(B) }2k\qquad \textbf{(C) }k\sqrt{k}\qquad \textbf{(D) }k^2\qquad  \textbf{(E) }y\sqrt{k}$

 \textbf{Answer: }D



\textbf{Solution}

We are given that 
 \[[f(x^2 + 1)]^{\sqrt(x)} = k\]We can rewrite $\frac{9+y^2}{y^2}$ as $\frac{9}{y^2} + 1$Thus, our function is now
 \[[f(\frac{9}{y^2} + 1)]^{\sqrt{\frac{12}{y}}} = k\]
 \[\Rrightarrow[f(\frac{9}{y^2} + 1)]^{\sqrt{\frac{3}{y} \cdot 4}} = k\]
 \[\Rrightarrow([f(\frac{9}{y^2} + 1)]^{\sqrt{\frac{3}{y}}})^{\sqrt{4}} = (k)^{\sqrt{4}}\]
 \[\Rrightarrow([f(\frac{9}{y^2} + 1)]^{\sqrt{\frac{3}{y}}})^{2} = (k)^{2} = k^2\]

 Thus, the answer is $k^2\implies \boxed{D}$.

 


	\item (\textit{AHSME 1978}) What is the smallest positive integer $n$ such that $\sqrt{n}-\sqrt{n-1}<.01$?


 $\textbf{(A) }2499\qquad \textbf{(B) }2500\qquad \textbf{(C) }2501\qquad \textbf{(D) }10,000\qquad  \textbf{(E) }\text{There is no such integer}$

 \textbf{Answer: }C



\textbf{Solution}

Adding $\sqrt{n - 1}$ to both sides, we get
 \[\sqrt{n} < \sqrt{n - 1} + 0.01.\]Squaring both sides, we get
 \[n < n - 1 + 0.02 \sqrt{n - 1} + 0.0001,\]which simplifies to
 \[0.9999 < 0.02 \sqrt{n - 1},\]or
 \[\sqrt{n - 1} > 49.995.\]Squaring both sides again, we get
 \[n - 1 > 2499.500025,\]so $n > 2500.500025$. The smallest positive integer $n$ that satisfies this inequality is $\boxed{2501}$.

 \textbf{alternative} Taking reciprocals and flipping the inequality we get 
 \[\sqrt{n}+\sqrt{n-1}>100\] Which is easy to see the answer is

 \begin{verbatim}
$\boxed{2501}$.
\end{verbatim}

-bjump

 


	\item (\textit{AHSME 1978}) A positive integer $n$ not exceeding $100$ is chosen in such a way that if $n\le 50$, then the probability of choosing 
$n$ is $p$, and if $n > 50$, then the probability of choosing $n$ is $3p$. The probability that a perfect square is chosen is


 $\textbf{(A) }.05\qquad \textbf{(B) }.065\qquad \textbf{(C) }.08\qquad \textbf{(D) }.09\qquad  \textbf{(E) }.1$

 \textbf{Answer: }C



\textbf{Solution}

Let's say that we will have $3$ slips for every number not exceeding $100$ but bigger than $50.$ This is to account for the $3p$ probability part. Let's now say that we will only have one slip for each number below or equal to $50.$ The probability(or $p$) will then be $\frac{1}{200}.$ Now let's have all the squares under $50,$ which are $1,4,9,16,25,36,49.$ The probability for these are $\frac{7}{200}.$ The numbers above $50$ that are squares are $64,81,100.$ We then need to multiply the probability by $3$ so the probability of these are $\frac{9}{200}.$ The answer is $\frac{7}{200}+\frac{9}{200}=0.08\implies\boxed{\textbf{(C).}}$



	\item (\textit{AHSME 1978}) If $a,b,c$ are non-zero real numbers such that 
 \[\frac{a+b-c}{c}=\frac{a-b+c}{b}=\frac{-a+b+c}{a},\] and 
 \[x=\frac{(a+b)(b+c)(c+a)}{abc},\] and $x<0,$ then $x$ equals


 $\textbf{(A) }{-}1\qquad \textbf{(B) }{-}2\qquad \textbf{(C) }{-}4\qquad \textbf{(D) }{-}6\qquad  \textbf{(E) }{-}8$

 \textbf{Answer: }A



\textbf{Solution}

From the equation 
 \[\frac{a+b-c}{c}=\frac{a-b+c}{b}=\frac{-a+b+c}{a},\] we add $2$ to each fraction to get 
 \[\frac{a+b+c}{c}=\frac{a+b+c}{b}=\frac{a+b+c}{a}.\]We perform casework on $a+b+c:$

            \begin{itemize}
        \item If $a+b+c\neq0,$ then $a=b=c,$ from which $x=\frac{(2a)(2a)(2a)}{a^3}=8.$ However, this contradicts the precondition $x<0.$
            \end{itemize}
            
            \begin{itemize}
        \item If $a+b+c=0,$ then $x=\frac{(-c)(-a)(-b)}{abc}=\boxed{\textbf{(A) }{-}1}.$
            \end{itemize}
            


	\item (\textit{AHSME 1978}) For all positive numbers $x$ distinct from $1$,


 
 \[\frac{1}{\log_3(x)}+\frac{1}{\log_4(x)}+\frac{1}{\log_5(x)}\]


 equals


 $\text{(A) }\frac{1}{\log_{60}(x)}\qquad\\ \text{(B) }\frac{1}{\log_{x}(60)}\qquad\\ \text{(C) }\frac{1}{(\log_{3}(x))(\log_{4}(x))(\log_{5}(x))}\qquad\\ \text{(D) }\frac{12}{\log_{3}(x)+\log_{4}(x)+\log_{5}(x)}\qquad\\ \text{(E) }\frac{\log_{2}(x)}{\log_{3}(x)\log_{5}(x)}+\frac{\log_{3}(x)}{\log_{2}(x)\log_{5}(x)}+\frac{\log_{5}(x)}{\log_{2}(x)\log_{3}(x)}$

 \textbf{Answer: }A



\textbf{Solution}


 \[\frac{1}{\log_3(x)}+\frac{1}{\log_4(x)}+\frac{1}{\log_5(x)}\implies \frac{1}{\frac{\log(x)}{\log(3)}}+\frac{1}{\frac{\log(x)}{\log(4)}}+\frac{1}{\frac{\log(x)}{\log(5)}}\implies \frac{\log(3)+\log(4)+\log(5)}{\log(x)}=\frac{\log(60)}{\log(x)}\implies \log_x(60)\implies \frac{1}{\log_{60}(x)}\]Thus, the answer is $(\textbf{A})$



	\item (\textit{AHSME 1978}) The following four statements, and only these are found on a card:

                    \begin{figure}[H]
                        \centering
                        \includegraphics[width=0.5\textwidth]{Problem/AHSME_I_1978_22_Problem_Figure_1.png}
                    \end{figure}
                    


 (Assume each statement is either true or false.) Among them the number of false statements is exactly


 $\textbf{(A)}\ 0 \qquad \textbf{(B)}\ 1 \qquad \textbf{(C)}\ 2 \qquad \textbf{(D)}\ 3 \qquad \textbf{(E)}\ 4$

 \textbf{Answer: }D



\textbf{Solution 1}

Any $2$ of these statements contradict each other, so at most $1$ statement is true.  Thus at least $3$ statements are false.  However, if all $4$ statements were false, then the last statement would be true, which would be a contradiction.  Therefore, there are $\boxed{\textbf{(D)}\ 3}$ false statements.  This is possible if only the third statement is true.  

 



\textbf{Solution 2}

If all of them are false, that would mean that the $4$th one is false too. Therefore, $E$ is not the correct answer. If exactly $3$ of them are false, that would mean that only $1$ statement is true. This is correct since if only $1$ statement is true, the card that is true is the one that has $3$ of these statements are false. If we have $1$ or $2$ false statements, that would mean that there is more than $1$ true statement. Therefore, our answer is $\boxed {(D)}$.

 


	\item (\textit{AHSME 1978}) 
                    \begin{figure}[H]
                        \centering
                        \includegraphics[width=0.5\textwidth]{Problem/AHSME_I_1978_23_Problem_Figure_1.png}
                    \end{figure}
                    


 Vertex $E$ of equilateral $\triangle ABE$ is in the interior of square $ABCD$, and $F$ is the point of intersection 
of diagonal $BD$ and line segment $AE$. If length $AB$ is $\sqrt{1+\sqrt{3}}$ then the area of $\triangle ABF$ is


 $\textbf{(A) }1\qquad \textbf{(B) }\frac{\sqrt{2}}{2}\qquad \textbf{(C) }\frac{\sqrt{3}}{2}\qquad \textbf{(D) }4-2\sqrt{3}\qquad  \textbf{(E) }\frac{1}{2}+\frac{\sqrt{3}}{4}$

 \textbf{Answer: }C



\textbf{Solution}

Place square $ABCD$ on the coordinate plane with $A$ at the origin. 

 In polar form, line $BD$ is $r\sin(\theta) = \sqrt{1 + \sqrt{3}} - r \cos(\theta)$ and line $AF$ is $\theta = \frac{\pi}{3}$.

 This means that the length from the origin to the intersection ($r$) is $r\frac{\sqrt{3}}{2} = \sqrt{1 + \sqrt{3}} - \frac{r}{2}$

 Solving for $r$, you get: $r = \frac{2}{\sqrt{1 + \sqrt{3}}}$

 Using the formula for area of a triangle ($A = \frac{ab\sin{C}}{2}$), you get $A = \frac{(\frac{2}{\sqrt{1 + \sqrt{3}}})(\sqrt{1 + \sqrt{3}}) \sin(\frac{\pi}{3})}{2} = \frac{\sqrt{3}}{2}$

 Getting $\textbf{C}$ as the answer

 


	\item (\textit{AHSME 1978}) If the distinct non-zero numbers $x ( y - z),~ y(z - x),~ z(x - y )$ form a geometric progression with common ratio $r$, 
then $r$ satisfies the equation


 $\textbf{(A) }r^2+r+1=0\qquad \textbf{(B) }r^2-r+1=0\qquad \textbf{(C) }r^4+r^2-1=0\qquad\\ \textbf{(D) }(r+1)^4+r=0\qquad  \textbf{(E) }(r-1)^4+r=0$

 \textbf{Answer: }A



\textbf{Solution}

Let the geometric progression be $a,$ $ar,$ $ar^2,$ so $a = x(y - z)$, $ar = y(z - x)$, and $ar^2 = z(x - y).$ Adding these equations, we get
 \[a + ar + ar^2 = xy - xz + yz - xy + xz - yz = 0.\]Since $a$ is nonzero, we can divide by $a$, to get $\boxed{r^2 + r + 1 = 0}$. The answer is (A).

 


	\item (\textit{AHSME 1978}) Let $a$ be a positive number. Consider the set $S$ of all points whose rectangular coordinates $(x, y )$ satisfy all of the following conditions:


 $\text{(i) }\frac{a}{2}\le x\le 2a\qquad \text{(ii) }\frac{a}{2}\le y\le 2a\qquad \text{(iii) }x+y\ge a\\ \\ \qquad \text{(iv) }x+a\ge y\qquad \text{(v) }y+a\ge x$


 The boundary of set S is a polygon with


 $\textbf{(A) }3\text{ sides}\qquad \textbf{(B) }4\text{ sides}\qquad \textbf{(C) }5\text{ sides}\qquad \textbf{(D) }6\text{ sides}\qquad \textbf{(E) }7\text{ sides}$

 \textbf{Answer: }D



\textbf{Solution}

Draw a picture. The answer is $\fbox{D}$.

 


	\item (\textit{AHSME 1978}) 
                    \begin{figure}[H]
                        \centering
                        \includegraphics[width=0.5\textwidth]{Problem/AHSME_I_1978_26_Problem_Figure_1.png}
                    \end{figure}
                    


 

In $\triangle ABC, AB = 10~ AC = 8$ and $BC = 6$. Circle $P$ is the circle with smallest radius which passes through $C$ 
and is tangent to $AB$. Let $Q$ and $R$ be the points of intersection, distinct from $C$ , of circle $P$ with sides $AC$ 
and $BC$, respectively. The length of segment $QR$ is


 $\textbf{(A) }4.75\qquad \textbf{(B) }4.8\qquad \textbf{(C) }5\qquad \textbf{(D) }4\sqrt{2}\qquad  \textbf{(E) }3\sqrt{3}$

 \textbf{Answer: }B



\textbf{Solution 1}

Let $D$ be the foot of the altitude from $C$ to $\overline{AB}$, and let $h = CD$.  Since the area of $\triangle ABC$ is $\frac{1}{2}\cdot 8\cdot 6= 24$, it follows that $\frac{1}{2}\cdot{10}\cdot{h} = 24$, so 

 


	\item (\textit{AHSME 1978}) There is more than one integer greater than $1$ which, when divided by any integer $k$ such that $2 \le k \le 11$, has a remainder of $1$. 
What is the difference between the two smallest such integers?


 $\textbf{(A) }2310\qquad \textbf{(B) }2311\qquad \textbf{(C) }27,720\qquad \textbf{(D) }27,721\qquad  \textbf{(E) }\text{none of these}$

 \textbf{Answer: }C



\textbf{Solution}

Let this integer be $n$. We have$n\equiv 1 \mod 2$, $n \equiv 1 \mod 3$, $n \equiv 1 \mod 4$ $\ldots$ $n \equiv 1 \mod 11$.Recall that if 
 \[a \equiv b \mod c\]and 
 \[a \equiv b \mod d\]then 
 \[a \equiv b \mod (lcm (c,d))\]We see that since$n\equiv 1 \mod 2$, $n \equiv 1 \mod 3$, $n \equiv 1 \mod 4$ $\ldots$ $n \equiv 1 \mod 11$.We have 
 \[n \equiv 1 \mod (lcm (2, 3, \ldots , 10, 11))\]

 From $2$ to $11$, $8$ contains the largest power of $2$, $9$ contains the largest power of $3$, and $10$ contains the largest power of $5$. Thus, our lcm is equal to
 \[2^3 \cdot 3^2 \cdot 5 \cdot 7 \cdot 11\]Since $n > 1$, our $2$ smallest values of $n$ are 
 \[n = 1 + (2^3 \cdot 3^2 \cdot 5 \cdot 7 \cdot 11)\] and
 \[n = 1 + 2(2^3 \cdot 3^2 \cdot 5 \cdot 7 \cdot 11)\]The difference between these values is simply the value of 
 \[(2^3 \cdot 3^2 \cdot 5 \cdot 7 \cdot 11) = 40 \cdot 99 \cdot 7\]
 \[= 3960 \cdot 7\]
 \[= 27720, \boxed{C}\]



	\item (\textit{AHSME 1978}) 
                    \begin{figure}[H]
                        \centering
                        \includegraphics[width=0.5\textwidth]{Problem/AHSME_I_1978_28_Problem_Figure_1.png}
                    \end{figure}
                    


 

If $\triangle A_1A_2A_3$ is equilateral and $A_{n+3}$ is the midpoint of line segment $A_nA_{n+1}$ for all positive integers $n$, 
then the measure of $\measuredangle A_{44}A_{45}A_{43}$ equals


 $\textbf{(A) }30^\circ\qquad \textbf{(B) }45^\circ\qquad \textbf{(C) }60^\circ\qquad \textbf{(D) }90^\circ\qquad  \textbf{(E) }120^\circ$

 \textbf{Answer: }E



\textbf{Solution (Repeating Cycles, Kinda )}

Upon graphing more points according to the rule mentioned in the problem, you can visibly see the rule repeats every 8 points (i.e. $A_9$ is the same as $A_1$, and $A_{10}$ is the same as $A_2$, etc.). Through very simple calculations, this means $A_{44}$ becomes $A_4$, $A_{45}$ becomes $A_5$, and $A_{43}$ becomes $A_3$, and you can clearly see on the diagram that this is a 120 degree angle, leading to $\fbox{E}$ as the final answer.

 


	\item (\textit{AHSME 1978}) Sides $AB,~ BC, ~CD$ and $DA$, respectively, of convex quadrilateral $ABCD$ are extended past $B,~ C ,~ D$ and $A$ to points $B',~C',~ D'$ and $A'$. 
Also, $AB = BB' = 6,~ BC = CC' = 7, ~CD = DD' = 8$ and $DA = AA' = 9$; and the area of $ABCD$ is $10$. The area of $A 'B 'C'D'$ is


 $\textbf{(A) }20\qquad \textbf{(B) }40\qquad \textbf{(C) }45\qquad \textbf{(D) }50\qquad  \textbf{(E) }60$

 \textbf{Answer: }D



\textbf{Solution}

Notice that the area of $\triangle$ $DAB$ is the same as that of $\triangle$ $A'AB$ (same base, same height). Thus, the area of $\triangle$ $A'AB$ is twice that (same height, twice the base). Similarly, [$\triangle$ $BB'C$] = 2 $\cdot$ [$\triangle$ $ABC$], and so on.

 Adding all of these, we see that the area the four triangles around $ABCD$ is twice [$\triangle$ $DAB$] + [$\triangle$ $ABC$] + [$\triangle$ $BCD$] + [$\triangle$ $CDA$], which is itself twice the area of the quadrilateral $ABCD$. Finally, [$A'B'C'D'$] = [$ABCD$] + 4 $\cdot$ [$ABCD$] = 5 $\cdot$ [$ABCD$] = $\fbox{50}$.

 Note: Anyone with a diagram would be of great help (still new to LaTex).

 


	\item (\textit{AHSME 1978}) In a tennis tournament, $n$ women and $2n$ men play, and each player plays exactly one match with every other player. 
If there are no ties and the ratio of the number of matches won by women to the number of matches won by men is $7/5$, then $n$ equals


 $\textbf{(A) }2\qquad \textbf{(B) }4\qquad \textbf{(C) }6\qquad \textbf{(D) }7\qquad  \textbf{(E) }\text{none of these}$

 \textbf{Answer: }E



\textbf{Solution 1}

Since there are $n$ women, the number of matches between only women is $\frac{n(n-1)}{2},$ and similarly, there are $(2n - 1)n$ matches between only men. Since every woman plays every man exactly once, there are $2n\cdot n = 2n^2$ matches which are between a man and a woman. Call these $2n^2$ matches co-ed matches, and let $w$ be the number of co-ed matches won by women.

 Then it follows that
 \[\frac 75 = \frac{\frac{n(n-1)}{2} + w}{(2n - 1)n + 2n^2 - w},\]which can be simplified to
 \[w = \frac{17}{8}n^2 - \frac38 n.\]

 The number of matches won by women must be less than the total number of matches, so we obtain the inequality
 \[\frac{17}{8}n^2 - \frac38 n\leq 2n^2.\]Rearranging and factoring gives
 \[0\leq -n(n-3),\]and the only integers which satisfy this inequality are $n = 0,1,2,$ and $3.$

 Clearly, there could not have been $0$ people in the tournament, so $n\neq 0.$ If $n = 1,$ then there would have been only one woman and two men in the tournament, in which case the woman could not have won the majority of the matches.

 We can now plug $n = 2$ back into the equation
 \[\frac 75 = \frac{\frac{n(n-1)}{2} + w}{(2n - 1)n + 2n^2 - w},\]and solving for $w$ gives $w = \frac{31}{4}.$ Since $w$ must be an integer, $n$ cannot be $2.$ It follows that $n = 3,$ so the answer is (E). (When $n = 3,$ solving gives $w = 18.$)

 



\textbf{Solution 2}

Since there are $n$ women and $2n$ men, there are a total of $n+2n=3n$ players. Hence, the number of total matches must be $\binom{3n}{2}=\frac{3n(3n-1)}{2}$. We also know that the ratio of the number of matches won by women to the number of matches won by men is $\frac{7}{5}$ and that there were no draws, so the total number of matches must be $(7+5)x=12x$ for some value $x$. This gives
 \[\frac{3n(3n-1)}{2}=12x\]So
 \[n(3n-1)=8x\]It follows that $n=8,4,3,2,1$ (various ways of splitting $8x$ into two factors). However, for solutions $n=4,2,1$, $x$ isn't an integer, therefore leaving solutions $n=8$ and $n=3$. 

 $n=8$ is invalid. This can be shown as when $n=8$,
 \[\frac{24(23)}{2}=12x\]
 \[x=23\]The ratio of matches won by gender is $\frac{7x}{5x}=\frac{161}{115}$, so the number of matches won by women must be at least $161$. The maximum number of matches won by women is equivalent to the number of matches between men (generates 1 match won by men no matter the outcome) subtracted from the total number of matches, which is $\binom{24}{2}-\binom{16}{2}=156$. As the maximum number of matches won is smaller than $161$, the solution is extraneous/invalid.

 Hence, the answer must be $n=3$, or $\boxed{\textbf{(E) }\text{none of these}}$




\textbf{Solution 3 (answer choices)}

Since $\frac{7}{12}$ of the total amount of games were won by women, the total amount of games has to be some multiple of $12$. The total amount of games is given by $\frac{(3n-1)(3n)}{2}$, when there are $3n$ total people. Plugging in $n=2$, we get $\frac{(5)(6)}{2}=15$, not a multiple of 12. Likewise for $n=4$, $\frac{(7)(8)}{2}=28$, which doesn't work either. Next, $n=6$ yields $\frac{(11)(12)}{2}=66$. Lastly, after testing out $n=7$ and obtaining $\frac{(13)(14)}{2}=91$, we see that the answer must be $\boxed{E}$.~martianrunner

 


\end{enumerate}